\clearpage
\VerdierPhysicalPageAnchor{0121}%
\begingroup
\small
\setlength{\emergencystretch}{1.5em}%
\setlength{\abovedisplayskip}{4pt plus 2pt minus 2pt}%
\setlength{\belowdisplayskip}{4pt plus 2pt minus 2pt}%

\VerdierSemanticAnchor{ch2-1-2-11-iv-to-i-complete}
prove the implication $\mathrm{iv)}\Longrightarrow\mathrm{i)}$. Thus let
$f\colon X\longrightarrow Y$ be a morphism equipped with a splitting
$g\colon Y\longrightarrow X$. The kernel functor of the morphism
$f\colon X\longrightarrow Y$ is then equal to the kernel functor of the
morphism $gf\colon X\longrightarrow X$, which is a projector of $X$
({\itshape i.e.} $gfgf=gf$).

\VerdierSemanticAnchor{ch2-1-2-11-projector-reduction}
We are therefore reduced to proving that, in a triangulated category in
which countable products are representable, projectors have kernels. Thus
let $e\colon X\longrightarrow X$ be a projector of $X$. Put $X_i=X$,
$i\in\mathbb N$, and let

\VerdierSemanticAnchor{ch2-1-2-11-phi}
\[
  \Phi\colon\prod_{i\in\mathbb N}X_i\longrightarrow
  \prod_{i\in\mathbb N}X_i, ,
\]
be the morphism defined by

\VerdierSemanticAnchor{ch2-1-2-11-1}
\begin{equation}
  \Phi(x_0,x_1,x_2,\ldots)
  =\bigl(ex_0+(1-e)x_1,ex_1+(1-e)x_2,\ldots\bigr)\,.
  \tag{1.2.11.1}
\end{equation}
Formula (1.2.11.1) defines a morphism by interpreting the symbols
$x_0,x_1,\ldots$ as morphisms from an arbitrary object of $\mathcal D$ to
the $X_i$. Likewise, let

\VerdierSemanticAnchor{ch2-1-2-11-psi}
\[
  \Psi\colon\prod_{i\in\mathbb N}X_i\longrightarrow
  \prod_{i\in\mathbb N}X_i, ,
\]
be the morphism defined by

\VerdierSemanticAnchor{ch2-1-2-11-2}
\begin{equation}
  \Psi(x_0,x_1,x_2,\ldots)
  =\bigl(x_0,(1-e)x_0+ex_1,(1-e)x_1+ex_2,\ldots\bigr)\,.
  \tag{1.2.11.2}
\end{equation}

\VerdierSemanticAnchor{ch2-1-2-11-kernel-conclusion}
It is clear that $\Phi\Psi=\operatorname{id}$ and hence that the morphism
$\Phi$ is an epimorphism. Therefore, by the implications already proved,
it has a representable kernel. On the other hand, the kernel functor of
$\Phi$ is canonically isomorphic to the kernel functor of
$e\colon X\longrightarrow X$. This functor is therefore representable and,
consequently, the projector $e$ has a kernel. This completes the proof of
Proposition (1.2.9).

\VerdierSemanticAnchor{ch2-1-2-12}
\medskip
\noindent\textbf{1.2.12.}\quad
Corollary (1.2.4) shows that the distinguished triangles containing a given
morphism are unique up to nonunique isomorphism.

\VerdierSemanticAnchor{ch2-1-2-12-functorial-question}
Moreover, axiom TRIII requires every morphism of morphisms to extend to a
morphism of triangles. One may therefore ask whether it is possible to
define the triangulated structure by functorially associating a
distinguished triangle with every morphism of the category. We shall show
that, if so, the category is decomposable under certain additional
hypotheses.

\VerdierSemanticAnchor{ch2-1-2-13}
\medskip
\noindent\textbf{Proposition 1.2.13.}\quad
\textit{Let $\mathcal D$ be a triangulated category and let
$\operatorname{Tr}$ be a functor from $\mathrm{Fl}_1(\mathcal D)$
(Chap.~I, 3.3.6), the category of morphisms of $\mathcal D$, to the
category of distinguished triangles (the full subcategory of the category
of triangles in $\mathcal D$ defined by the distinguished triangles):}

\VerdierSemanticAnchor{ch2-1-2-13-tr-assignment}
\[
  \operatorname{Tr}\!\left(X\xrightarrow{f}Y\right)
  =X\xrightarrow{f}Y\xrightarrow{a(f)}\tau(f)
  \xrightarrow{b(f)}X[1], .
\]
\endgroup
